% =====================================================================
%  Project Proposal
%  BangHallu: A Bengali Corpus for Hallucination Detection in LLMs
%
%  Compile with:  pdflatex proposal.tex   (run twice for the ToC)
%
%  This file is deliberately pure ASCII. pdflatex with inputenc utf8 cannot
%  typeset Bengali script -- it stops with "Unicode character not set up
%  for use with LaTeX" -- so the single Bengali title in the reference list
%  is given in romanised form. To use real Bengali script, compile with
%  XeLaTeX or LuaLaTeX and load a Bengali-capable font, for example:
%      \usepackage{fontspec}
%      \newfontfamily\bengali{Noto Serif Bengali}[Script=Bengali]
%  Every package used below ships with a standard TeX Live install.
% =====================================================================

\documentclass[12pt,a4paper]{article}

\usepackage[utf8]{inputenc}
\usepackage[T1]{fontenc}
\usepackage[margin=1in]{geometry}
\usepackage{times}
\usepackage{setspace}
\usepackage{booktabs}
\usepackage{array}
\usepackage{amsmath}
\usepackage{enumitem}
\usepackage{fancyhdr}
\usepackage[hidelinks]{hyperref}
\usepackage{titlesec}

\onehalfspacing
\setlength{\parskip}{6pt}
\setlength{\parindent}{0pt}

\titleformat{\section}{\large\bfseries}{\thesection.}{0.6em}{}
\titlespacing*{\section}{0pt}{14pt}{6pt}
\titleformat{\subsection}{\normalsize\bfseries}{\thesubsection}{0.6em}{}

\pagestyle{fancy}
\fancyhf{}
\fancyhead[L]{\small Hallucination Detection for Bengali QA}
\fancyhead[R]{\small KUET}
\fancyfoot[C]{\thepage}
\renewcommand{\headrulewidth}{0.4pt}

\newcolumntype{L}[1]{>{\raggedright\arraybackslash}p{#1}}

\begin{document}

% =====================================================================
%  TITLE PAGE
% =====================================================================
\begin{titlepage}
\centering
\vspace*{1.2cm}

{\LARGE\bfseries BangHallu: A Bengali Corpus for Hallucination Detection in Large Language Models\par}

\vspace{1.6cm}
{\large\bfseries Project Proposal\par}

\vspace{2.0cm}

{\large\bfseries Submitted by\par}
\vspace{0.4cm}
{\large
Tawhidul Hasan\\
2107004\\[0.35cm]
Md. Saif Ahmed Shejan\\
2107009\par}

\vspace{1.6cm}

{\large\bfseries Submitted to\par}
\vspace{0.4cm}
{\large
Dr. K. M. Azharul Hasan\\
Professor\\[0.35cm]
Md. Nazirulhasan Shawon\\
Assistant Professor\par}

\vspace{2.0cm}

{\large Department of Computer Science and Engineering\par}
{\large\bfseries Khulna University of Engineering \& Technology\par}

\vfill
\end{titlepage}

\newpage
\tableofcontents
\newpage

% =====================================================================
\section{Abstract}
% =====================================================================

Large language models now answer questions in Bengali fluently. They are also
sometimes simply wrong, and they are wrong in a confident voice.

We propose to build a detector that reads a question, an optional supporting
passage, and a candidate answer, and decides whether that answer is
correct or hallucinated. It will not try to produce the right
answer, only to judge the answer it is shown.

The project delivers two things: a Bengali question-answering corpus of reasonable pairs where each pair being one question with one correct and one
hallucinated answer, and a benchmark of models climbing from simple
word-counting methods through embeddings and recurrent networks to pretrained
encoders such as BanglaBERT. This proposal is the outcome of a preparatory study of the Bengali and
multilingual hallucination literature, of candidate data sources, and of several
possible designs.

% =====================================================================
\section{Problem Statement}
% =====================================================================

\subsection{The practical problem}

Bengali is spoken by roughly 230 million people and is the sixth most widely
spoken language in the world, yet it remains low-resource for evaluation: there
are far fewer test sets, benchmarks and trained models for Bengali than for
English. Students across Bangladesh increasingly use language models as study
aids for competitive and school examinations. The danger is specific. When a model does not know an answer, it rarely says so.
It produces something that reads like an answer, but the fact inside it is
invented.

\subsection{Why this is hard to detect}

Our preparatory study identified several difficulties the design must address.

\begin{enumerate}[leftmargin=*,itemsep=3pt]
  \item Wrong answers look like right ones. A hallucinated answer
  reuses the correct vocabulary and grammar, often differing by a single year,
  name, or character. Surface similarity offers little signal.

  \item Two different judgements are needed. With a passage, the
  question is ``does this passage support this answer?'', that is reading
  comprehension. Without one, it becomes ``is this true?'', that is factual recall.

  \item Datasets can contain shortcuts. If correct answers are copied
  word-for-word from the passage while wrong answers are not, a model can score
  highly by matching strings and learn nothing about meaning.

\end{enumerate}

\subsection{The gap this project fills}

BenHalluEval (2026) is the first dedicated Bengali hallucination benchmark, but
it evaluates how well existing models detect hallucination by prompting
them. To our knowledge no published work trains a dedicated detector for
Bengali and compares model families on a like-for-like basis. That is the gap: a
trained Bengali detector, benchmarked across a full ladder of methods, on a
corpus whose validity is explicitly audited rather than assumed.

% =====================================================================
\section{Objectives}
% =====================================================================

\subsection{Primary objectives}

\begin{enumerate}[label=O\arabic*.,leftmargin=*,itemsep=3pt]
  \item Build a validated Bengali hallucination corpus of QA
  pairs, covering both grounded  and closed-book conditions, labelled by native Bengali speakers.

  \item Establish a reliable annotation protocol, demonstrated by two annotators
  independently labelling a blind sample and reaching Cohen's $\kappa \geq 0.60$
  before full annotation begins.

  \item Benchmark a complete model ladder, from lexical n-gram models through
  embeddings and recurrent networks to pretrained transformer encoders.

  \item Reach macro-F1 $\geq 0.80$ on the has-context hard subset, the
  pairs a string matcher cannot solve and $\geq 0.60$ on no-context.
  
\end{enumerate}

\subsection{Secondary objectives}

Establish Bengali baselines against which a second phase, covering romanised Bengali (Banglish), can later be measured.


% =====================================================================
\section{Literature Review}
% =====================================================================

\subsection{Hallucination, taxonomy, and shortcut learning}

Hallucination is conventionally split in two. Intrinsic hallucination
contradicts a source document supplied to the model; extrinsic
hallucination cannot be checked against any supplied source and is a claim about
the world. This maps directly onto our two conditions, which is why our taxonomy
has one branch for each: \texttt{entity}, \texttt{numeric}, \texttt{relational}
and \texttt{contradiction} for grounded items; \texttt{fabricated} and
\texttt{overclaim} for closed-book ones.

A recurring failure in dataset construction, characterised by Geirhos et al.
(2020), is that a model learns a surface regularity instead of the intended
task. In our setting the obvious candidate is string overlap between answer and
passage. The standard remedy is a probe: train a deliberately impoverished
classifier and check that it fails.

\subsection{Bengali resources}

BenHalluEval (arXiv:2605.31483, 2026) is the first benchmark built for
Bengali hallucination: four tasks, 12{,}000 hallucinated candidates generated
across twelve error types, nine models evaluated under a dual-track protocol.
Three of its findings shape our design. Dual-track evaluation is
necessary, measuring only detection rate hides a model that calls everything
hallucinated, so we report per-class metrics separately. Native-speaker
validation reached $\kappa = 0.911$--$0.926$, which puts our $\kappa \geq 0.60$
requirement in perspective as a floor, not a target. Code-mixed input
behaves differently from native script, which is why romanised Bengali is a
separate phase rather than a variant.

Several encoders cover Bengali and form the upper rungs of our ladder.
BanglaBERT is an ELECTRA discriminator pretrained on
Bengali and our primary candidate.

\subsection{Methodological templates}

BanTH (arXiv:2410.13281) is our closest model: binary classification
over transliterated Bengali, with a published further-pretraining recipe and
annotation guidelines whose structure we intend to imitate. Its most useful
lesson is sobering, across seven very different encoders macro-F1 spanned only
77.35 to 74.51, a range of 2.8 points, and further-pretraining gains were around
two points, slightly negative for XLM-R. Architecture choice alone will not
transform a result; input format and further pretraining matter more.

ViHallu (arXiv:2601.04711), a Vietnamese shared task with 111 teams,
reached 84.80 macro-F1 at best against an encoder-only baseline of 32.83.
MedHallu contributes two ideas we adopt: stratifying a corpus by
difficulty, and letting annotators say ``not sure'' instead of guessing.

% =====================================================================
\section{Dataset and Data Collection}
% =====================================================================

\subsection{Sources and construction}

Two source families have been identified: Bengali Wikipedia,
supplying passage text for grounded items, and BCS-style question
banks, supplying chiefly the closed-book items. Each question will be paired with two answers: one correct and one hallucinated.
The hallucinated answer will be produced from the same source text with
language-model assistance, under prompt constraints on length and register so
the two are not separable by superficial features. Every stage will be a
deterministic script at a fixed seed, so the corpus rebuilds byte-for-byte.


Subjects will span reading comprehension, history, grammar, mathematics, idiom,
antonym, synonym, vocabulary, law, science, BCS general knowledge and
literature. A per-subject cap is needed because mathematics alone dominates the
closed-book pool. No subject will be excluded for being difficult.

Splits are grouped by pair so both answers to a question land together.
Splitting by record would put the correct answer in training and the wrong one
in test, letting a model score by recognising the question rather than judging
the answer. The test split will be examined exactly once, at the end; every
intermediate decision is made on development data.

\subsection{Annotation protocol}

Two native speakers will annotate independently, following guidelines written
before annotation begins. For each item they record the binary label, the
error type when the answer is wrong, a difficulty judgement, and a note.

Two rules govern outside knowledge. For no-context items an annotator
may verify a fact and should note that they did. For has-context items they must
not: the question is whether this passage supports the answer, and an
external source cannot settle that. Where an item cannot be labelled honestly,
the passage does not address the question, the answer type does not match, or
the text is broken, the annotator writes \texttt{unsure} rather than guessing.
Such items are excluded from the agreement score and reported by subject, since
a cluster in one subject is a finding about the corpus, not a failure by the
annotator.

Agreement is a gate. A blind sample of 100 items will be labelled by
both annotators and scored with Cohen's $\kappa$. If $\kappa < 0.60$ we treat the
guidelines as the problem rather than the annotators, write an explicit rule for
each disagreement, and repeat on the same items. Only then does full annotation
begin. Effort is then allocated by importance: the test split labelled in full
by both with disagreements adjudicated, development by one annotator, and
training through a 20\% sample sufficient to estimate label noise.

% =====================================================================
\section{Methodology}
\label{sec:methodology}
% =====================================================================

\subsection{Task formulation and input}

Given a question $q$, an optional passage $c$, and a candidate answer $a$, we
learn $f(q, c, a) \rightarrow \{0, 1\}$, where $1$ means correct and
$0$ means hallucinated. Text will be normalised with the \texttt{csebuetnlp} normaliser before
BanglaBERT, which was pretrained on normalised text. Three input formats will be
compared, since format frequently matters more than architecture: F1
(question + answer, the only option without a passage), F2 (context +
question + answer).

Truncation needs care. At a fixed maximum length a long passage is cut,
and if the cut removes the supporting sentence the label is no longer derivable
from the input and the model trains on noise. We will truncate the
context only, never the question or the answer.

\subsection{The model ladder}

Models are built in ascending order of capacity, so each rung is an
interpretable floor for the next.

\begin{center}
\begin{tabular}{L{2.4cm} L{5.4cm} L{5.4cm}}
\toprule
\textbf{Rung} & \textbf{Models} & \textbf{Purpose} \\
\midrule
Lexical & $n$-gram + LogReg / SVM; TF-IDF over word 1--2 and character 3--5
grams & Lower bound. Character features should matter given Bengali inflection \\
\addlinespace
Embedding & Skip-gram / word2vec + LogReg; + XGBoost &
Is distributional meaning alone enough? \\
\addlinespace
Recurrent & BiRNN; BiLSTM; BiLSTM + attention &
Do word order and selective attention help? \\
\addlinespace
Pretrained & BanglaBERT, MuRIL, XLM-R, mBERT, IndicBERT v2 &
The real contenders \\
\addlinespace
Adapted & Further pretraining on unlabelled Bengali (mBERT, XLM-R only) &
Domain adaptation; BanglaBERT is ELECTRA and cannot use masked-LM \\
\addlinespace
Ensemble & Soft-voting over the strongest three encoders &
Typically worth one to three points \\
\addlinespace
Reference & Zero-shot / few-shot LLM on 300--500 items &
A ceiling marker, not a competitor \\
\bottomrule
\end{tabular}
\end{center}

Classical models are expected to reach roughly 0.50-0.65 macro-F1 at this
corpus size. That will be the correct outcome, not a failure; they have
no pretraining to draw on, and their limits are what motivate the upper rungs.

Hyperparameters start from BanTH's published setup: learning rate
$2 \times 10^{-5}$, AdamW, max length 256, batch 32, 5 epochs with early
stopping, warmup 0.1, weight decay 0.01, fp16. Every configuration runs with
three seeds (42, 1337, 2024), reported as mean $\pm$ standard deviation, with
five-fold cross-validation over train and dev while test stays untouched.


% =====================================================================
\section{Evaluation Metrics}
% =====================================================================

\subsection{Metric and breakdowns}

The primary metric is macro-F1, the unweighted mean of both classes' F1
scores, chosen because it weights correct and hallucinated equally so a model
cannot score well by favouring the majority class. For a class $c$:
\[
F1_c = \frac{2 P_c R_c}{P_c + R_c},
\qquad
\text{macro-F1} = \tfrac{1}{2}\left(F1_{\text{correct}} + F1_{\text{halluc.}}\right).
\]
We also report accuracy, per-class precision and recall, AUROC, and the
confusion matrix.

A single averaged number conceals exactly the differences that matter, so every
result is decomposed four ways: has-context vs no-context (two
different tasks); easy vs hard (without it a string matcher looks
competitive); by subject (law, science and BCS need outside knowledge a
closed-book model lacks, and a collapse there is a finding, not a defect); and
by hallucination type, once annotation is complete.

\subsection{Targets and gates}

\begin{center}
\begin{tabular}{L{6.2cm} c L{4.4cm}}
\toprule
\textbf{Metric} & \textbf{Target} & \textbf{Role} \\
\midrule
Macro-F1, has-context hard subset & $\geq 0.80$ & headline result \\
Macro-F1, has-context overall & report only & shown beside the string baseline \\
Macro-F1, no-context & $\geq 0.60$ & second headline \\
Metadata shortcut probe & $< 0.60$ & blocking gate \\
Cohen's $\kappa$ & $\geq 0.60$ & gate before full annotation \\
\bottomrule
\end{tabular}
\end{center}

% =====================================================================
\section{Expected Outcomes}
% =====================================================================

\subsection{Deliverables}

An annotated Bengali hallucination corpus of at least 4{,}000 pairs with locked,
reproducible splits; a written annotation guideline document revised against
real disagreements; an inter-annotator agreement report; a source and licence
register; trained model checkpoints across the ladder; five results tables (main
results, difficulty breakdown, per-type F1, format ablation, shortcut audit); a
complete experiment log including failures; a validity audit report; a
reproducible codebase; and a technical report.

\subsection{Expected performance}

Calibrated from comparable published work rather than chosen optimistically.

\begin{center}
\begin{tabular}{L{5.2cm} c L{4.8cm}}
\toprule
\textbf{Setting} & \textbf{Realistic macro-F1} & \textbf{Trigger for audit} \\
\midrule
Has-context (grounded) & 0.80 -- 0.90 & $> 0.95$ \\
No-context (closed book) & 0.60 -- 0.75 & $> 0.85$ \\
Hard subset only & 0.55 -- 0.70 & --- \\
Classical models & 0.50 -- 0.65 & expected, not a failure \\
Zero-shot LLM reference & 0.60 -- 0.72 & --- \\
\bottomrule
\end{tabular}
\end{center}

\subsection{Findings we anticipate}

\begin{itemize}[leftmargin=*,itemsep=3pt]
  \item No-context will score well below has-context. Closed-book
  verification without retrieval is genuinely hard. That is a correct result and
  we will report it plainly.
  \item The easy-hard gap will be large. It may be the single most
  informative number the project produces, since it separates models that read
  from models that pattern-match.
  \item Model spread will be narrow. BanTH found 2.8 points across
  seven encoders; we will report ties honestly rather than manufacture a ranking.
  \item Models may collapse on law, science and BCS. Reported per
  subject, that is a finding about the limits of the approach.
\end{itemize}

% =====================================================================
\section{Discussion}
\label{sec:discussion}
% =====================================================================

\subsection{Conclusions reached during exploration}

Four decisions were not obvious at the outset, and each rules out an approach
that might otherwise seem natural.

1. Bengali script before romanised Bengali. Romanised Bengali was our original
starting point, but it is a harder problem layered on an unsolved one, with no
Bengali baseline to measure degradation against. Our exploration of automatic
transliteration also showed that rule-based character mapping produces text a
native speaker struggles to read, which would make reliable annotation
impossible. It becomes a second phase.

2. Measure difficulty rather than filter it. We considered excluding law and
science, reasoning that an annotator who must look a fact up cannot label
reliably. We concluded against it: removing difficult material makes a benchmark
easier without making it better, and conceals a limitation instead of measuring
it. A difficulty field and a per-subject breakdown replace exclusion.

3. Question answering only. A cloze item trains a model to copy a missing span
rather than judge whether an answer is supported, and is trivially solved by
string matching. This concerns task type, not difficulty.

4. No retrieval at inference. Retrieval would let a system look the fact up,
measuring the retriever rather than the detector, and would empty the
closed-book condition of meaning.

\subsection{Why the shortcut audit matters most}

The most likely way for this project to fail is not a low score but a high one
that means nothing. If correct answers turn out to be verbatim spans of their
passages while hallucinated ones are not, a rule containing no learning could
score highly, and a model reporting a similar figure would look like a success
while representing almost no improvement. The whole evaluation design is
arranged around detecting that: the metadata probe is a blocking gate, and the
hard subset rather than the overall figure is the headline target. The same
reasoning explains why pairs whose incorrect answer is also a verbatim span are
kept rather than discarded; those are the ones a model cannot cheat on.

\subsection{Anticipated limitations}

\begin{itemize}[leftmargin=*,itemsep=2pt]
  \item At roughly 4{,}000 pairs the corpus is modest, a realistic constraint of
  a low-resource language.
  \item Hallucinated answers will be produced with language-model assistance, so
  a model could learn the generator's fingerprint instead of the concept. The
  answer-only probe is our defence.
  \item The system says an answer is wrong; it does not say which words are
  wrong, nor produce a correction.
  \item Agreement between two annotators is less robust than across three.
\end{itemize}

\subsection{Risks and mitigations}

\begin{center}
\begin{tabular}{L{5.0cm} L{8.6cm}}
\toprule
\textbf{Risk} & \textbf{Mitigation} \\
\midrule
A shortcut inflates scores & Metadata probe as a blocking gate, run on the pilot
and on every rebuild \\
Agreement falls below 0.60 & Fix the guidelines, not the annotators; rerun the
same items \\
Question-bank licence forbids redistribution & Settle provenance before
ingestion; keep those items separable \\
Model differences too small to distinguish & Three seeds, McNemar's test,
bootstrap CIs; report ties \\
No-context performance is poor & Expected; report as a finding, not a defect \\
Test set contaminated by repeated use & All decisions on dev; test evaluated once \\
\bottomrule
\end{tabular}
\end{center}

% =====================================================================
\section{Conclusion}
% =====================================================================

A language model answering a Bengali examination question can be confidently
wrong, and the student least able to detect the error is exactly the student who
asked. We propose to build a detector that judges whether a given answer is
correct or hallucinated, together with the corpus needed to train and evaluate
it fairly.

The plan follows from a preparatory study rather than assumption: we chose
Bengali script over romanised Bengali, chose to measure difficulty rather than
filter it out, ruled out fill-in-the-blank items as a different task, and ruled
out retrieval as something that would measure the wrong thing. Two commitments
follow from that. We will measure what a rule with no learning in it can score
on our own data before trusting any model, and report the hard subset as the
headline result. And subjects that are hard to verify will stay in the corpus
and be measured, because a benchmark made easier is not a benchmark made better.

If the work succeeds, the outcome will be the first trained Bengali
hallucination detector with a published model ladder and an audited corpus.

% =====================================================================
\section*{References}
\addcontentsline{toc}{section}{References}
% =====================================================================

\begin{enumerate}[leftmargin=*,itemsep=3pt,label={[\arabic*]}]

\item BenHalluEval: A Multi-Task Hallucination Evaluation Framework for Large
Language Models on Bengali. arXiv:2605.31483, 2026.

\item BanTH: A Multi-label Hate Speech Detection Dataset for Transliterated
Bangla. arXiv:2410.13281, 2024.

\item DSC2025 ViHallu: Vietnamese Hallucination Detection Shared Task.
arXiv:2601.04711, 2026.

\item Bhattacharjee, A., Hasan, T., Ahmad, W. U., and Shahriyar, R. BanglaBERT:
Language Model Pretraining and Benchmarks for Low-Resource Language
Understanding Evaluation in Bangla. arXiv:2101.00204, 2021.

\item Khanuja, S. et al. MuRIL: Multilingual Representations for Indian
Languages. arXiv:2103.10730, 2021.

\item Conneau, A. et al. Unsupervised Cross-lingual Representation Learning at
Scale (XLM-R). ACL, 2020.

\item Devlin, J., Chang, M.-W., Lee, K., and Toutanova, K. BERT: Pre-training of
Deep Bidirectional Transformers for Language Understanding. NAACL, 2019.

\item Clark, K., Luong, M.-T., Le, Q. V., and Manning, C. D. ELECTRA:
Pre-training Text Encoders as Discriminators Rather Than Generators. ICLR, 2020.

\item Doddapaneni, S. et al. IndicBERT v2: Towards Leaving No Indic Language
Behind. ACL, 2023.

\item MedHallu: A Comprehensive Benchmark for Detecting Medical Hallucinations
in Large Language Models. \texttt{medhallu.github.io}, 2025.

\item Mikolov, T., Chen, K., Corrado, G., and Dean, J. Efficient Estimation of
Word Representations in Vector Space. ICLR Workshop, 2013.

\item Hochreiter, S., and Schmidhuber, J. Long Short-Term Memory. Neural
Computation, 9(8), 1997.

\item Cohen, J. A Coefficient of Agreement for Nominal Scales. Educational and
Psychological Measurement, 20(1), 1960.

\item McNemar, Q. Note on the Sampling Error of the Difference Between
Correlated Proportions or Percentages. Psychometrika, 12(2), 1947.

\item Geirhos, R. et al. Shortcut Learning in Deep Neural Networks. Nature
Machine Intelligence, 2020.

\end{enumerate}

\end{document}
